
En este capítulo se presentan los objetivos y el catálogo de requisitos del nuevo sistema informático. Opcionalmente, se incluye el análisis de la brecha entre los requisitos planteados y la solución base seleccionada. 

\section{Objetivos de Negocio}
Esta sección debe contener los objetivos de negocio que se esperan alcanzar cuando el sistema software a desarrollar esté en producción. Adicionalmente, se pueden incluir los modelos de procesos, que normalmente son los modelos de procesos de negocio actuales con ciertas mejoras.

\section{Objetivos del Sistema}
Esta sección debe contener la especificación de los objetivos o requisitos generales del sistema.

\section{Requisitos funcionales}
Descripción completa de la funcionalidad que ofrece el sistema. 

\section{Requisitos no funcionales}
Descripción de otros requisitos (relacionados con la calidad del software) que el sistema deberá satisfacer: portabilidad, seguridad, mantenibilidad, accesibilidad, usabilidad, fiabilidad, eficiencia, etc. Puede usarse de referencia lo indicado en la norma ISO/IEC 25010.

\section{Requisitos de información}
En esta sección se describen los requisitos de gestión de información (datos) que el sistema debe gestionar.

\section{Reglas de negocio}
En el desarrollo del sistema, hay que tener en cuenta las denominadas reglas de negocio, es decir, el conjunto de restricciones, normas o políticas de la organización que deben ser respetadas por el sistema, las cuales suelen ser cambiantes.

\section{Análisis GAP}
Esta sección es opcional y es de aplicación cuando en la revisión de la técnica se hayan localizado soluciones informáticas que permitan servir de base para el nuevo sistema software. En esta sección se debe identificar y medir las diferencias entre lo que proporcionan dichas soluciones y los requisitos definidos para el proyecto. El resultado de este análisis permitirá identificar cuáles de éstos requisitos ya están solventados total o parcialmente por los distintos sistemas y poder elegir uno de ellos. 